\chapter{Zastosowania rachunku całkowego}
\label{ma1:ch:integral-applications}

Całka Riemanna pozwala przejść od skończonych przybliżeń geometrycznych do dokładnych wielkości. W tym rozdziale pracujemy z krzywymi i obszarami o regularności wystarczającej do takich przybliżeń. Nie używamy ogólnej teorii miary ani ogólnej definicji pola powierzchni rozmaitości.

\section{Pola obszarów płaskich}
\begin{twierdzenie}
\label{ma1:thm:area-between-graphs}
Niech $a<b$, a $u,v:[a,b]\to\R$ będą ciągłe i $u\le v$. Pole obszaru
\[
 D=\{(x,y):a\le x\le b,\ u(x)\le y\le v(x)\}
\]
wynosi
\[
 A(D)=\int_a^b(v(x)-u(x))\,\dd x.
\]
\end{twierdzenie}
\begin{proof}
Pole rozumiemy w sensie Jordana: przybliżamy obszar od wewnątrz i od zewnątrz skończonymi sumami prostokątów o rozłącznych wnętrzach. Jeżeli pola takich przybliżeń można uczynić dowolnie bliskimi, ich wspólna wartość graniczna jest polem (każde przybliżenie wewnętrzne ma pole nie większe niż dowolne zewnętrzne).

Dla podziału $P$ przedziału $[a,b]$ oznaczmy jego podprzedziały przez $J_i$, ich długości przez $\Delta x_i$, a kresy funkcji $u,v$ na $J_i$ odpowiednio przez $m_{u,i},M_{u,i},m_{v,i},M_{v,i}$. Prostokąt zewnętrzny na $J_i$ ma wysokość $M_{v,i}-m_{u,i}$, a wewnętrzny wysokość $\max(0,m_{v,i}-M_{u,i})$; gdy drugi przedział wysokości jest pusty, pomijamy prostokąt. Różnica wysokości jest nie większa niż
\[
 (M_{v,i}-m_{v,i})+(M_{u,i}-m_{u,i}).
\]
Istotnie, zastąpienie $\max(0,m_{v,i}-M_{u,i})$ przez $m_{v,i}-M_{u,i}$ może tylko zwiększyć różnicę. Z jednostajnej ciągłości obu funkcji (wniosek \ref{ma1:cor:interval-uniform}) dla dowolnego $\varepsilon>0$ wybieramy średnicę podziału tak małą, aby każda oscylacja była mniejsza od $\varepsilon/(2(b-a))$. Wówczas różnica pól przybliżeń jest mniejsza od $\varepsilon$.
Dla każdego punktu wyróżnionego $\xi_i\in J_i$ wysokość $v(\xi_i)-u(\xi_i)$ leży między tymi dwiema wysokościami. Sumy Riemanna dla $v-u$ są więc ściśnięte między polami przybliżeń. Z twierdzenia \ref{ma1:thm:riemann-darboux-equivalence} ich granica wynosi $\int_a^b(v-u)$, co kończy dowód.
\end{proof}
Jeżeli dwa wykresy $f,g$ przecinają się, przyjmujemy $v=\max(f,g)$, $u=\min(f,g)$. Są to funkcje ciągłe, ponieważ $\max(f,g)=(f+g+|f-g|)/2$ i $\min(f,g)=(f+g-|f-g|)/2$. Pole między wykresami wynosi więc $\int_a^b|f-g|$. Całka bez modułu może mieć znak ujemny lub zerowy i nie musi być polem.

\begin{przyklad}
\label{ma1:ex:area-parabola}
Dla $0\le x\le1$ mamy $x^2\le x$. Pole między wykresami wynosi
\[
 \int_0^1(x-x^2)\,\dd x=\frac12-\frac13=\frac16.
\]
\end{przyklad}
\begin{center}
\begin{tikzpicture}[x=5cm,y=3.5cm,>=Stealth]
 \fill[blue!15] (0,0)--(1,1)--plot[domain=1:0,samples=70] (\x,{\x*\x})--cycle;
 \draw[->] (-.08,0)--(1.15,0) node[right] {$x$};
 \draw[->] (0,-.08)--(0,1.13) node[above] {$y$};
 \draw[blue!70!black,thick] (0,0)--(1,1) node[pos=.75,above left] {$y=x$};
 \draw[teal!70!black,thick,domain=0:1,samples=70] plot (\x,{\x*\x});
 \node[teal!70!black,right] at (.79,.61) {$y=x^2$};
 \fill[orange!30] (.48,{.55*.55}) rectangle (.55,.48);
 \draw[orange!80!black] (.48,{.55*.55}) rectangle (.55,.48);
 \draw[<->] (.60,{.55*.55})--(.60,.48) node[midway,right] {$v-u$};
 \node[below] at (1,0) {$1$};
\end{tikzpicture}

\small Rysunek: pole między wykresami i przykładowy prostokąt wewnętrzny. Jego wysokość przybliża różnicę wartości funkcji na krótkim przedziale.
\end{center}

\section{Długość krzywej}
\begin{definicja}
\label{ma1:def:curve-length}
Dla ciągłej krzywej parametrycznej $\gamma:[a,b]\to\R^d$ i podziału $P=(t_0,\ldots,t_n)$ definiujemy długość łamanej wpisanej
\[
 L_P(\gamma)=\sum_{i=1}^n\|\gamma(t_i)-\gamma(t_{i-1})\|_2.
\]
Długość krzywej to $L(\gamma)=\sup_P L_P(\gamma)$, z dopuszczeniem $+\infty$. Krzywą o skończonej długości nazywamy prostowalną. Jest to długość drogi z uwzględnieniem wielokrotnych przejść, a niekoniecznie długość samego obrazu.
\end{definicja}
\begin{twierdzenie}
\label{ma1:thm:curve-length}
Jeżeli $\gamma=(\gamma_1,\ldots,\gamma_d)$ jest klasy $C^1$ na $[a,b]$ (pochodne na końcach rozumiemy jednostronnie), to
\[
 L(\gamma)=\int_a^b\|\gamma'(t)\|_2\,\dd t
 =\int_a^b\sqrt{\sum_{j=1}^d\gamma_j'(t)^2}\,\dd t.
\]
Nie wymagamy, aby $\gamma'(t)$ była niezerowa.
\end{twierdzenie}
\begin{proof}
Dla $u<v$ z Newtona--Leibniza w każdej współrzędnej mamy $\gamma(v)-\gamma(u)=\int_u^v\gamma'(s)\,\dd s$, gdzie całka wektorowa oznacza tu całkę każdej współrzędnej. Jeżeli przyrost $w=\gamma(v)-\gamma(u)$ jest niezerowy, bierzemy jednostkowy wektor $e=w/\|w\|_2$. Nierówność Cauchy'ego--Schwarza w $\R^d$ daje
\[
 \|w\|_2=e\cdot w=\int_u^v e\cdot\gamma'(s)\,\dd s
 \le\int_u^v\|\gamma'(s)\|_2\,\dd s.
\]
Dla $w=0$ nierówność jest oczywista. Sumując po podprzedziałach, otrzymujemy $L_P\le\int_a^b\|\gamma'\|_2$ dla każdego $P$.

Dla przeciwnego oszacowania niech $\varepsilon>0$. Każda współrzędna $\gamma'$ jest jednostajnie ciągła, więc istnieje $\delta>0$ takie, że $|s-t|<\delta$ pociąga $\|\gamma'(s)-\gamma'(t)\|_2<\varepsilon$ (wystarczy oszacować każdą współrzędną przez $\varepsilon/\sqrt d$). Dla przedziału $[u,v]$ o długości mniejszej od $\delta$ i $\xi\in[u,v]$ stosujemy poprzednie oszacowanie normy całki do ciągłej funkcji wektorowej $\gamma'(s)-\gamma'(\xi)$: identyczny argument z wektorem jednostkowym daje
\[
 \|\gamma(v)-\gamma(u)-(v-u)\gamma'(\xi)\|_2
 \le\varepsilon(v-u).
\]
Odwrotna nierówność trójkąta daje
\[
 \big|\|\gamma(v)-\gamma(u)\|_2-(v-u)\|\gamma'(\xi)\|_2\big|
 \le\varepsilon(v-u).
\]
Po zsumowaniu błąd między długością łamanej a sumą Riemanna funkcji $\|\gamma'\|_2$ nie przekracza $\varepsilon(b-a)$. Wybierając podziały o średnicach dążących do zera, otrzymujemy zbieżność długości łamanych do wskazanej całki. Ponieważ każda z tych długości jest nie większa niż supremum, całka jest nie większa niż $L(\gamma)$. Oba oszacowania dają równość.
\end{proof}
Dla wykresu $y=f(x)$ klasy $C^1$ stosujemy $\gamma(x)=(x,f(x))$, zatem
\[
 L=\int_a^b\sqrt{1+f'(x)^2}\,\dd x.
\]
Jeśli krzywa składa się ze skończenie wielu części klasy $C^1$, sumujemy ich długości: rozdrobnienie podziału nie zmniejsza długości łamanej dzięki nierówności trójkąta, więc można ograniczyć supremum do podziałów zawierających wszystkie punkty sklejenia.

\begin{przyklad}
\label{ma1:ex:circle-length}
Dla $\gamma(t)=(r\cos t,r\sin t)$, $0\le t\le2\pi$, $r>0$, mamy $\|\gamma'(t)\|_2=r$. Długość jednego pełnego obiegu okręgu wynosi $2\pi r$. Dla $0\le t\le4\pi$ wynosi $4\pi r$, mimo identycznego obrazu krzywej.
\end{przyklad}
\begin{wniosek}
\label{ma1:cor:arc-reparametrization}
Jeśli $\gamma:[a,b]\to\R^d$ jest klasy $C^1$, a $\varphi:[c,d]\to[a,b]$ jest klasy $C^1$, monotoniczna i surjektywna, to $L(\gamma\circ\varphi)=L(\gamma)$.
\end{wniosek}
\begin{proof}
Reguła łańcuchowa w każdej współrzędnej daje $\|(\gamma\circ\varphi)'(s)\|_2=\|\gamma'(\varphi(s))\|_2|\varphi'(s)|$. Gdy $\varphi$ jest niemalejąca, $\varphi'\ge0$, a jej wartości na końcach to $a,b$. Podstawienie (twierdzenie \ref{ma1:thm:integration-substitution}) daje szukaną równość; funkcję $\|\gamma'\|_2$ można ciągle przedłużyć poza $[a,b]$ wartościami stałymi. Gdy $\varphi$ jest nierosnąca, $\varphi'\le0$ i końce mają wartości $b,a$; minus z modułu znosi minus z orientacji całki.
\end{proof}

\section{Objętości brył obrotowych}
\begin{twierdzenie}
\label{ma1:thm:rotation-volume}
Niech $u,v:[a,b]\to\R$ będą ciągłe i $0\le u\le v$. Obracając obszar między ich wykresami wokół osi $x$, otrzymujemy bryłę
\[
 B=\{(x,y,z):a\le x\le b,\ u(x)^2\le y^2+z^2\le v(x)^2\}.
\]
Jej objętość wynosi
\[
 V(B)=\pi\int_a^b(v(x)^2-u(x)^2)\,\dd x.
\]
\end{twierdzenie}
\begin{proof}
Pole koła o promieniu $r$ wynosi $\pi r^2$: dla $r>0$ stosujemy twierdzenie \ref{ma1:thm:area-between-graphs} do wykresów $\pm\sqrt{r^2-x^2}$ na $[-r,r]$ i przykład \ref{ma1:ex:trig-root-substitution}; dla $r=0$ pole jest zerowe. Pole pierścienia o promieniach $0\le p\le q$ wynosi $\pi(q^2-p^2)$ przez odjęcie pól kół (okrąg ma pole zerowe: każdy z dwóch wykresów ciągłych można pokryć prostokątami o sumie pól dowolnie małej, stosując twierdzenie \ref{ma1:thm:area-between-graphs} z równymi funkcjami ograniczającymi). Objętość walca o takim przekroju i długości $h$ wynosi $h\pi(q^2-p^2)$: każde wewnętrzne i zewnętrzne prostokątne przybliżenie przekroju mnożymy przez odcinek długości $h$, otrzymując odpowiednie przybliżenia prostopadłościanami.

Na podprzedziale $J_i$ podziału $[a,b]$ bryła mieści się w walcu pierścieniowym o promieniach $\min_{J_i}u$ i $\max_{J_i}v$. Zawiera walec o promieniach $\max_{J_i}u$ i $\min_{J_i}v$, jeżeli pierwszy nie przekracza drugiego; w przeciwnym przypadku bierzemy pusty zbiór. Różnica pól przekrojów jest ograniczona przez $\pi(\operatorname{osc}_{J_i}u^2+\operatorname{osc}_{J_i}v^2)$, gdzie oscylacja oznacza różnicę kresu górnego i dolnego. Dowód oszacowania jest taki jak dla prostokątów w twierdzeniu \ref{ma1:thm:area-between-graphs}, zastosowanym do $u^2,v^2$. Z jednostajnej ciągłości kwadratów i sumowania długości podprzedziałów różnica objętości dąży do zera przy średnicy podziału dążącej do zera. Walce można z kolei przybliżyć skończonymi sumami prostopadłościanów z dowolnie małym łącznym błędem. Otrzymujemy istnienie objętości w sensie Jordana. Suma Riemanna funkcji $\pi(v^2-u^2)$ leży między tymi objętościami i dąży do całki, co dowodzi wzoru.
\end{proof}
W szczególności dla $u=0$ otrzymujemy metodę krążków, a dla $u>0$ metodę pierścieni. Nie wolno bez dodatkowego rozpatrzenia stosować różnicy kwadratów do obszaru przecinającego oś obrotu, gdyż jego części mogą po obrocie się pokrywać.

\begin{center}
\begin{tikzpicture}[x=1.15cm,y=1cm,>=Stealth]
 \draw[->] (-.4,0)--(5,0) node[right] {$x$};
 \draw[thick,blue!70!black] (0,1.2)--(4,1.7);
 \draw[thick,blue!70!black] (0,-1.2)--(4,-1.7);
 \draw[thick,teal!70!black] (0,.45)--(4,.7);
 \draw[thick,teal!70!black] (0,-.45)--(4,-.7);
 \draw[dashed] (0,0) ellipse (.26 and 1.2);
 \draw[blue!70!black] (4,0) ellipse (.34 and 1.7);
 \draw[teal!70!black] (4,0) ellipse (.14 and .7);
 \fill[blue!22,even odd rule] (2,0) ellipse (.29 and 1.45) (2,0) ellipse (.12 and .575);
 \draw[blue!70!black] (2,0) ellipse (.29 and 1.45);
 \draw[teal!70!black] (2,0) ellipse (.12 and .575);
 \draw[->] (2,0)--(2,1.45) node[midway,left] {$v(x)$};
 \draw[->] (2,0)--(2,-.575) node[midway,right] {$u(x)$};
\end{tikzpicture}

\small Rysunek: przekrój bryły płaszczyzną prostopadłą do osi obrotu jest pierścieniem o polu $\pi(v^2-u^2)$.
\end{center}
\begin{przyklad}
\label{ma1:ex:sphere-volume}
Kula promienia $r>0$ powstaje przez obrót półkola $v(x)=\sqrt{r^2-x^2}$, $u=0$, $-r\le x\le r$. Zatem
\[
 V=\pi\int_{-r}^r(r^2-x^2)\,\dd x
 =\pi[r^2x-x^3/3]_{-r}^r=\frac43\pi r^3.
\]
\end{przyklad}

\section{Pole powierzchni obrotowej}
\begin{twierdzenie}
\label{ma1:thm:rotation-surface}
Niech $\gamma(t)=(x(t),y(t))$ będzie klasy $C^1$ na $[a,b]$, $y\ge0$, a $x$ ściśle monotoniczna. Pole powierzchni powstałej przez obrót krzywej wokół osi $x$, w sensie granicy pól bocznych stożków ściętych wyznaczonych przez wpisane łamane, wynosi
\[
 S=2\pi\int_a^b y(t)\sqrt{x'(t)^2+y'(t)^2}\,\dd t.
\]
\end{twierdzenie}
\begin{proof}
Dla kolejnych punktów łamanej o promieniach $p,q\ge0$ i długości odcinka $\ell$ pole boczne stożka ściętego wynosi $\pi(p+q)\ell$. Dla $p=q$ jest to pole boczne walca (obwód podstawy razy długość tworzącej), także gdy oba promienie są zerowe. Dla $p\ne q$ rozwinięcie powierzchni bocznej jest różnicą dwóch współśrodkowych sektorów: ich promienie tworzących $l_1<l_2$ spełniają $l_2-l_1=\ell$ i długości łuków wynoszą $2\pi\min(p,q)$ i $2\pi\max(p,q)$. Pole sektora to połowa iloczynu długości łuku i promienia (uzasadnienie wzoru sektorowego podajemy w następnej sekcji). Różnica pól wynosi $\pi(\max(p,q)l_2-\min(p,q)l_1)$. Z podobieństwa trójkątów $l_1=\ell\min(p,q)/|p-q|$, $l_2=\ell\max(p,q)/|p-q|$, więc po podstawieniu otrzymujemy $\pi(p+q)\ell$. Przypadek promienia zerowego odpowiada pełnemu stożkowi.

Dla podziału $P$ pole boczne przybliżenia wynosi
\[
 S_P=\pi\sum_i(y(t_{i-1})+y(t_i))\|\gamma(t_i)-\gamma(t_{i-1})\|_2.
\]
Ścisła monotoniczność $x$ zapobiega nakładaniu się wnętrz kolejnych powierzchni bocznych. Niech $M=\max y$, $B=\max\|\gamma'\|_2$. Dla $\varepsilon>0$ i dostatecznie małej średnicy podziału jednostajna ciągłość daje $|y(t_i)-y(\xi_i)|,|y(t_{i-1})-y(\xi_i)|\le\varepsilon$ oraz, jak w dowodzie twierdzenia \ref{ma1:thm:curve-length},
\[
 \big|\|\gamma(t_i)-\gamma(t_{i-1})\|_2-\|\gamma'(\xi_i)\|_2\Delta t_i\big|
 \le\varepsilon\Delta t_i.
\]
Ponadto długość każdego odcinka łamanej nie przekracza $B\Delta t_i$. Odejmując od $S_P$ sumę Riemanna funkcji $2\pi y\|\gamma'\|_2$, najpierw zastępujemy sumę promieni przez $2y(\xi_i)$, a potem długość odcinka przez $\|\gamma'(\xi_i)\|_2\Delta t_i$. Błąd jest nie większy niż $2\pi\varepsilon(B+M)(b-a)$. Sumy Riemanna dążą do całki, więc wszystkie takie przybliżenia mają tę samą granicę. Jest to pole w podanym sensie.
\end{proof}
Dla wykresu $y=f(x)\ge0$ klasy $C^1$ otrzymujemy $S=2\pi\int_a^b f(x)\sqrt{1+f'(x)^2}\,\dd x$. Wzór dotyczy powierzchni bocznej, nie dodaje pól podstaw.

\begin{przyklad}
\label{ma1:ex:sphere-surface}
Półokrąg parametryzujemy jako $(r\cos t,r\sin t)$ dla $0\le t\le\pi$. Pierwsza współrzędna jest ściśle malejąca, druga nieujemna, a prędkość wynosi $r$. Pole sfery wynosi
\[
 S=2\pi r^2\int_0^\pi\sin t\,\dd t=4\pi r^2.
\]
Parametryzacja pozwala uniknąć nieograniczonej pochodnej wykresu półokręgu na końcach.
\end{przyklad}

\section{Współrzędne biegunowe}
Dla $r\ge0$ i kąta $\theta$ zapisujemy $(x,y)=(r\cos\theta,r\sin\theta)$. Aby uniknąć wielokrotnego pokrycia obszaru, w twierdzeniu o polu przyjmiemy przedział kątów długości nie większej od $2\pi$.

Pole sektora koła promienia $r$ i kąta $\omega\in[0,2\pi]$ wynosi $r^2\omega/2$. Istotnie, dzielimy kąt na $n$ równych części $\delta=\omega/n<\pi$. W każdej wpisujemy trójkąt o dwóch bokach długości $r$ i kącie $\delta$, o polu $r^2\sin\delta/2$. Zewnętrzny trójkąt ograniczony tymi samymi półprostymi i styczną prostopadłą do dwusiecznej ma podstawę $2r\tan(\delta/2)$ i wysokość $r$, więc pole $r^2\tan(\delta/2)$. Dla $r>0$, $\omega>0$ każda z sum pól (wewnętrznych i zewnętrznych) dąży do $r^2\omega/2$, gdyż $\sin\delta/\delta\to1$ i $\tan(\delta/2)/(\delta/2)\to1$ (granice elementarne z rozdziału \ref{ma1:ch:elementary-functions}). Dla $\omega=0$ lub $r=0$ pole jest zerowe. Przybliżenia trójkątne można przybliżyć prostokątami z dowolnie małym błędem: trójkąt po podziale wzdłuż pionowych prostych przechodzących przez wierzchołki jest sumą obszarów między funkcjami afinicznymi, a twierdzenie \ref{ma1:thm:area-between-graphs} stosuje się do każdej części. Długość łuku sektora wynosi $r\omega$ z twierdzenia \ref{ma1:thm:curve-length}, co uzasadnia także wzór na pole sektora jako połowę iloczynu promienia i długości łuku.

\begin{twierdzenie}
\label{ma1:thm:polar-area}
Niech $\alpha<\beta\le\alpha+2\pi$ oraz $\rho:[\alpha,\beta]\to[0,\infty)$ będzie ciągła. Pole obszaru
\[
 D=\{(r\cos\theta,r\sin\theta):\alpha\le\theta\le\beta,\ 0\le r\le\rho(\theta)\}
\]
wynosi $A=\tfrac12\int_\alpha^\beta\rho(\theta)^2\,\dd\theta$.
\end{twierdzenie}
\begin{proof}
Na każdej części podziału kątów $J_i$ obszar zawiera sektor o promieniu $\min_{J_i}\rho$ i mieści się w sektorze o promieniu $\max_{J_i}\rho$. Wnętrza sektorów dla różnych przedziałów kątów są rozłączne; wspólne półproste mają pole zerowe. Różnica sum pól zewnętrznych i wewnętrznych jest równa połowie różnicy sum Darboux funkcji $\rho^2$. Dąży do zera dzięki ciągłości $\rho^2$. Przybliżenia sektorami można zastąpić prostokątnymi z dowolnie małym błędem, więc pole istnieje i jest równe połowie całki.
\end{proof}
\begin{wniosek}
\label{ma1:cor:polar-length}
Dla $\rho$ klasy $C^1$ długość krzywej $\gamma(\theta)=(\rho(\theta)\cos\theta,\rho(\theta)\sin\theta)$ wynosi
\[
 L=\int_\alpha^\beta\sqrt{\rho'(\theta)^2+\rho(\theta)^2}\,\dd\theta.
\]
\end{wniosek}
\begin{proof}
Pochodne współrzędnych to $\rho'\cos\theta-\rho\sin\theta$ oraz $\rho'\sin\theta+\rho\cos\theta$. Po podniesieniu do kwadratu i zsumowaniu wyrazy mieszane znoszą się, a $\sin^2\theta+\cos^2\theta=1$. Stosujemy twierdzenie \ref{ma1:thm:curve-length}.
\end{proof}
\begin{przyklad}
\label{ma1:ex:cardioid-area}
Dla kardioidy $\rho(\theta)=1+\cos\theta$, $0\le\theta\le2\pi$, mamy
\[
 A=\frac12\int_0^{2\pi}(1+2\cos\theta+\cos^2\theta)\,\dd\theta
 =\frac12(2\pi+0+\pi)=\frac{3\pi}{2},
\]
gdzie użyliśmy $\cos^2\theta=(1+\cos2\theta)/2$.
\end{przyklad}
\begin{center}
\begin{tikzpicture}[scale=1.55,>=Stealth]
 \fill[blue!12] plot[domain=0:360,samples=181,smooth] ({(1+cos(\x))*cos(\x)},{(1+cos(\x))*sin(\x)})--cycle;
 \draw[->] (-.65,0)--(2.35,0) node[right] {$x$};
 \draw[->] (0,-1.5)--(0,1.5) node[above] {$y$};
 \fill[orange!35] (0,0)--plot[domain=38:55,samples=25] ({(1+cos(\x))*cos(\x)},{(1+cos(\x))*sin(\x)})--cycle;
 \draw[blue!70!black,thick] plot[domain=0:360,samples=181,smooth] ({(1+cos(\x))*cos(\x)},{(1+cos(\x))*sin(\x)});
 \draw[->] (0,0)--({(1+cos(45))*cos(45)},{(1+cos(45))*sin(45)}) node[midway,above left] {$\rho(\theta)$};
 \draw (.34,0) arc[start angle=0,end angle=45,radius=.34];
 \node at (.5,.2) {$\theta$};
\end{tikzpicture}

\small Rysunek: kardioida i wąski wycinek kątowy, którego pole przybliżamy polem sektora koła. Kąty we wzorach całkowych mierzymy w radianach.
\end{center}

\section{Środki ciężkości i momenty bezwładności}
Rozważmy jednorodną płytę $D=\{(x,y):a\le x\le b,\ u(x)\le y\le v(x)\}$ o ciągłych brzegach i dodatnim polu $A$. Przyjmujemy jednostkową gęstość powierzchniową, więc masa jest równa $A$. Moment pierwszego rzędu względem osi $x$ to suma graniczna iloczynów mas małych elementów przez ich współrzędne $y$ ze znakiem; analogicznie moment względem osi $y$ wykorzystuje współrzędną $x$. Moment bezwładności względem osi jest sumą graniczną iloczynów mas przez kwadraty odległości od osi.

\begin{twierdzenie}
\label{ma1:thm:lamina-moments}
Dla opisanej płyty środek ciężkości ma współrzędne
\[
 \bar x=\frac1A\int_a^b x(v-u)\,\dd x,\qquad
 \bar y=\frac1{2A}\int_a^b(v^2-u^2)\,\dd x.
\]
Momenty bezwładności względem osi $x,y$ wynoszą odpowiednio
\[
 I_x=\frac13\int_a^b(v^3-u^3)\,\dd x,\qquad
 I_y=\int_a^b x^2(v-u)\,\dd x.
\]
\end{twierdzenie}
\begin{proof}
Dla ciągłej funkcji wagi $w(x,y)$ na prostokącie zawierającym $D$ wykażemy, że suma graniczna ważonych pól jest równa
\[
 \int_a^b\left(\int_{u(x)}^{v(x)}w(x,y)\,\dd y\right)\dd x.
\]
Nie jest to założenie ogólnego twierdzenia Fubiniego. Weźmy prostokątne przybliżenia wewnętrzne i zewnętrzne z dowodu twierdzenia \ref{ma1:thm:area-between-graphs}, a następnie rozdrobnijmy je także w kierunku pionowym, aby średnice prostokątów dążyły do zera. Ponieważ $w$ jest ograniczona, wkład obszaru pomiędzy przybliżeniami ma moduł nie większy niż $\max|w|$ razy różnica pól, a więc dąży do zera. Z jednostajnej ciągłości $w$ (na zwartym prostokącie) wynika, że zmiana punktów wyróżnionych w prostokątach zmienia sumę o najwyżej maksymalną oscylację na prostokątach razy pole prostokąta zawierającego $D$. Zatem ten błąd także dąży do zera.

Na pionowym pasku $J_i$ zastępujemy $w(x,y)$ przez $w(\xi_i,y)$ dla $\xi_i\in J_i$. Jednostajna ciągłość daje błąd nie większy niż $\varepsilon$ razy pole paska, gdy średnica podziału jest dostatecznie mała. Zastąpienie pionowych końców przez $u(\xi_i),v(\xi_i)$ zmienia ważoną sumę o najwyżej $\max|w|(\operatorname{osc}_{J_i}u+\operatorname{osc}_{J_i}v)\Delta x_i$. Sumując, otrzymujemy błąd dążący do zera. W każdym tak otrzymanym prostokątnym pasku sumy po pionowych podziałach dążą do $\Delta x_i\int_{u(\xi_i)}^{v(\xi_i)}w(\xi_i,y)\,\dd y$ z całkowalności funkcji ciągłej jednej zmiennej. Funkcja $x\mapsto\int_{u(x)}^{v(x)}w(x,y)\,\dd y$ jest ciągła: różnicę jej wartości szacujemy przez długość wspólnego przedziału ograniczającą błąd zmiany wagi oraz $\max|w|$ razy sumę zmian końców, analogicznie do dowodu reguły Leibniza \ref{ma1:thm:parameter-leibniz}, bez używania pochodnych. Stąd sumy po paskach dążą do podanej całki Riemanna.

Stosujemy ten wynik kolejno do $w=x,y,y^2,x^2$. Całkowanie po $y$ daje odpowiednio $x(v-u)$, $(v^2-u^2)/2$, $(v^3-u^3)/3$, $x^2(v-u)$. Środek ciężkości definiujemy przez podzielenie momentów pierwszego rzędu przez masę, co jest możliwe, bo $A>0$. Otrzymujemy wszystkie cztery wzory.
\end{proof}
\begin{przyklad}
\label{ma1:ex:triangle-centroid}
Dla trójkąta $0\le x\le1$, $0\le y\le x$ mamy $A=1/2$. Zatem
\[
 \bar x=2\int_0^1 x^2\,\dd x=\frac23,\qquad
 \bar y=\int_0^1 x^2\,\dd x=\frac13,
 \qquad I_x=\frac1{12},\quad I_y=\frac14.
\]
Ostatnie dwie wartości wynikają z całek $\frac13\int_0^1 x^3\,\dd x$ oraz $\int_0^1 x^3\,\dd x$.
\end{przyklad}

\section{Twierdzenia Pappusa--Guldina}
\begin{twierdzenie}
\label{ma1:thm:pappus-volume}
Niech $D$ będzie jednorodnym obszarem z twierdzenia \ref{ma1:thm:lamina-moments}, z $0\le u\le v$ i $A>0$. Objętość bryły powstałej przez obrót wokół osi $x$ wynosi
\[
 V=2\pi\bar y A.
\]
\end{twierdzenie}
\begin{proof}
Z twierdzenia \ref{ma1:thm:rotation-volume} i wzoru na moment pierwszego rzędu mamy $V=\pi\int_a^b(v^2-u^2)=2\pi A\bar y$. Czynnik $2\pi\bar y$ jest długością drogi środka ciężkości w czasie jednego pełnego obrotu.
\end{proof}
\begin{twierdzenie}
\label{ma1:thm:pappus-surface}
Niech $\gamma(t)=(x(t),y(t))$ spełnia założenia twierdzenia \ref{ma1:thm:rotation-surface} i ma dodatnią długość $L$. Dla jednorodnego drutu o jednostkowej gęstości liniowej jego środek ciężkości ma drugą współrzędną
\[
 \bar y=\frac1L\int_a^b y(t)\|\gamma'(t)\|_2\,\dd t.
\]
Pole powierzchni obrotowej wynosi $S=2\pi\bar y L$.
\end{twierdzenie}
\begin{proof}
Moment drutu jest granicą sum $\sum_i y(\xi_i)\|\gamma(t_i)-\gamma(t_{i-1})\|_2$. Z oszacowania długości odcinka z dowodu twierdzenia \ref{ma1:thm:curve-length} różnica między tą sumą a sumą Riemanna funkcji $y\|\gamma'\|_2$ jest nie większa niż $\max|y|\varepsilon(b-a)$ dla dostatecznie drobnego podziału. Granica jest więc równa podanej całce; podzielenie przez masę $L>0$ daje $\bar y$. Wzór na $S$ wynika z twierdzenia \ref{ma1:thm:rotation-surface}.
\end{proof}
Twierdzenia sformułowaliśmy dla obszarów i krzywych leżących po jednej stronie osi oraz dla powierzchni o jednoznacznych przekrojach wzdłuż osi. Ogólniejsze wersje wymagają osobnej kontroli wielokrotności pokrycia.

\begin{przyklad}
\label{ma1:ex:torus-volume}
Niech $R>r>0$. Obrót koła o promieniu $r$ i środku $(0,R)$ wokół osi $x$ daje pełny torus. Na $[-r,r]$ mamy $u=R-\sqrt{r^2-x^2}>0$, $v=R+\sqrt{r^2-x^2}$. Ponieważ $u+v=2R$, wzór na środek ciężkości daje $\bar y=R$. Pole koła wynosi $\pi r^2$, więc objętość torusa to $2\pi^2Rr^2$. Pole powierzchni można obliczyć, rozdzielając okrąg na dwie części monotoniczne względem $x$ i parametryzując obie przez $(r\cos t,R+r\sin t)$ na $[0,\pi]$ i $[\pi,2\pi]$. Powierzchnie obu części mają rozłączne wnętrza, a ich wspólne okręgi na końcach nie wnoszą pola bocznego. Z twierdzenia \ref{ma1:thm:rotation-surface}
\[
 S=2\pi r\int_0^{2\pi}(R+r\sin t)\,\dd t=4\pi^2Rr.
\]
\end{przyklad}

W każdym zastosowaniu najpierw określamy obiekt geometryczny, następnie wielkości elementarnych przybliżeń i ich granicę. Wzory z całkami nie usuwają potrzeby sprawdzenia regularności, znaków promieni oraz tego, czy parametryzacja nie pokrywa wielokrotnie tego samego obiektu.
